\documentclass[10pt]{article}
\usepackage{icmlko}

\icmlmeta{IP 파운데이션 월드모델}{322}{검사는 거르지 줄 세우지 못한다}{전용 축이 값을 하는 조건을 예측했다가 반증됐다 --- 그리고 대신 나온 것도 표본 셋짜리다}

\begin{document}

\twocolumn[
\icmlheader

\begin{abstract}
노트 321이 남긴 물음: \texttt{fund\_cat} 은 펀딩 $+$0.117 을 버는데
\texttt{mkt\_cat} 은 시장팝업 $+$0.026 이다. 검사 ①②③④ 를 둘 다 통과하고
\texttt{mkt\_cat} 은 \textbf{유보에서도 라벨을 가른다}($H{=}14.4$,
$p{=}0.026$). 유일하게 다른 것이 \textbf{도메인 몫}으로 보였다 --- 펀딩
3.08\% 대 시장팝업 0.97\%, 노트 318의 절벽을 사이에 두고. 가설을 세웠다:
\textbf{전용 축은 그 도메인이 절벽 위에 있어야 값을 한다.} 노트 319 · 320의
가중으로 시장팝업 몫을 0.97\% $\to$ 2.5\% $\to$ 9.09\% 로 올려 다시 쟀다.
\textbf{반증됐다} --- \texttt{mkt\_cat} 이득이 $+$0.0259 $\to$ $-$0.0058
$\to$ $-$0.0233 으로 \textbf{몫을 올릴수록 나빠진다}(예측은 $+$0.05$\sim$
$+$0.20 이었다). \texttt{fund\_cat} 은 $+$0.117 $\to$ $+$0.131 $\to$
$+$0.080 으로 버틴다. \textbf{그러면 무엇이 다른가} --- 다섯 성질을 재
봤다. \textbf{유보 효과 크기는 오히려 거꾸로다}(\texttt{fund\_cat} $\epsilon^2
{=}0.037$ 로 제일 낮은데 이득 최고, \texttt{anime\_medium} 0.112 로 제일
높은데 음수). 세 축을 맞게 정렬하는 것은 \textbf{학습 구간 $\epsilon^2$}
(0.148 $>$ 0.127 $>$ 0.005)와 \textbf{기준 $\rho$}(0.254 $<$ 0.410 $<$
0.476, 거꾸로) 둘뿐이다. \textbf{그런데 $n{=}3$ 이다} --- 노트 315 · 316의
잣대로 세 값의 순서는 주장할 수 없다. 그래서 \textbf{다음 축을 지을 때
시험할 예측으로 등록만 한다.} \textbf{검사 넷은 거르기용이고, 무엇이 벌지는
지어 봐야 안다.}
\end{abstract}
]

\section{가설과 반증}

노트 321이 끝나면서 남긴 자리 --- 검사 넷을 다 통과한 축 둘이 스무 배
차이 난다.

\begin{center}\footnotesize
\begin{tabular}{lrrl}
\hline
& 도메인 몫 & 유보 이득 & \\
\hline
\texttt{fund\_cat} & 3.08\% & $+$0.1173 & 절벽 안쪽 \\
\texttt{mkt\_cat} & 0.97\% & $+$0.0259 & 절벽 아래 \\
\hline
\end{tabular}
\end{center}

\texttt{mkt\_cat} 은 축이 나쁜 게 아니다 --- 학습 $H{=}18.4$($p{=}0.010$) ·
\textbf{유보 $H{=}14.4$($p{=}0.026$)} 로 유보에서도 라벨을 가른다.
\textbf{모형이 안 쓴다.} 노트 318의 논리를 전용 축에 옮기면 가설이 나온다:
\textbf{0.97\% 짜리 도메인의 전용 축으로 가르는 분할은 탐욕 기준에서 안
뽑힌다.}

\paragraph{예측을 적고 쟀다} \texttt{preregister\_privateweight.json} 에
``$\alpha{=}0$ 에서 \textbf{크게 오른다} --- $+$0.05$\sim$$+$0.20''을 적고,
반증 조건도 함께 적었다: ``안 오르면 몫이 아니라 다른 것이다.''

\begin{figure*}[t]
\centering
\includegraphics[width=\textwidth]{figs/rank.pdf}
\caption{왼쪽 --- 가중으로 몫을 올리며 잰 전용 축의 이득. 오른쪽 --- 다섯
성질 중 어느 것이 세 축을 정렬하나.}
\end{figure*}

\begin{center}\small
\begin{tabular}{lrrr}
\hline
& $\alpha{=}1$ & $\alpha{=}0.75$ & $\alpha{=}0$ \\
시장팝업 몫 & 0.97\% & 2.5\% & 9.09\% \\
\hline
\textbf{\texttt{mkt\_cat}} & $+$0.0259 & \textbf{$-$0.0058} & \textbf{$-$0.0233} \\
\texttt{fund\_cat} & $+$0.1173 & $+$0.1311 & $+$0.0799 \\
\hline
\end{tabular}
\end{center}

\textbf{몫을 올릴수록 나빠진다.} 예측과 방향이 반대다. \textbf{가설은
죽는다.}

\section{그러면 무엇이 다른가}

다섯 성질을 재고 이득 순서(\texttt{fund} $>$ \texttt{mkt} $>$
\texttt{anime})와 맞는지 봤다.

\begin{center}\footnotesize
\begin{tabular}{lrrrl}
\hline
& \texttt{fund} & \texttt{mkt} & \texttt{anime} & \\
\hline
\textbf{유보 이득} & $+$0.117 & $+$0.026 & $-$0.009 & --- \\
\hline
학습 $\epsilon^2$ & 0.148 & 0.127 & 0.005 & \textbf{맞음} \\
기준 $\rho$ & 0.254 & 0.410 & 0.476 & \textbf{거꾸로} \\
\hline
유보 $\epsilon^2$ & 0.037 & 0.089 & 0.112 & 거꾸로(틀림) \\
최빈 밖 & 77\% & 80\% & 6\% & 못 맞춤 \\
도메인 몫 & 3.1\% & 1.0\% & 14.1\% & 못 맞춤 \\
\hline
\end{tabular}
\end{center}

\paragraph{유보 효과 크기는 거꾸로다} \texttt{fund\_cat} 이 유보에서
제일 약하게 가르는데($\epsilon^2{=}0.037$, $p{=}0.24$) 이득이 제일 크고,
\texttt{anime\_medium} 이 제일 세게 가르는데($\epsilon^2{=}0.112$,
$H{=}69.6$) 음수다. \textbf{축이 유보를 잘 가른다고 모형이 잘 쓰는 게
아니다.}

\paragraph{맞는 둘} \textbf{학습 $\epsilon^2$}(효과 크기, $p$ 값이 아니다)와
\textbf{기준 $\rho$}(낮을수록 벌 여지가 크다)가 세 축을 맞게 정렬한다.
둘 다 \textbf{유보를 안 보고 계산된다} --- 쓸 수 있는 자리다.

\section{그런데 $n{=}3$ 이다}

\begin{center}\small
\begin{tabular}{ll}
\hline
세 값의 순서가 우연히 맞을 확률 & $1/6$ \\
뒤집혀 맞을 것까지 세면 & $1/3$ \\
\hline
\end{tabular}
\end{center}

\textbf{노트 315 · 316이 정한 잣대로는 주장할 수 없다} --- 오늘 아침에
``못 가르는 값들의 순서도 못 가른다''로 노트 하나를 물렀는데, 여기서 표본
셋짜리 순서로 법칙을 세우면 같은 잘못이다.

\paragraph{그래서 등록만 한다} 다음에 전용 축을 지을 때 \textbf{학습
$\epsilon^2$ 와 기준 $\rho$ 를 먼저 재고 이득을 예측한 뒤} 짓는다. 맞으면
검사 다섯째가 되고 틀리면 오늘의 셋이 우연이었던 것이다.

\section{그리고 검사의 자리를 다시 적는다}

\begin{center}\footnotesize
\begin{tabular}{ll}
\hline
검사가 하는 것 & \\
\hline
① 정보가 없는 축 & 거른다 \\
② 시기에 흔들리는 축 & 거른다 \\
③ 이미 있는 축 & 거른다 \\
④ 거의 상수인 축 & 거른다 \\
\hline
\textbf{무엇이 벌지} & \textbf{못 고른다} \\
\hline
\end{tabular}
\end{center}

넷을 다 통과한 축 셋이 $+$0.117 · $+$0.026 · $-$0.009 다. \textbf{검사는
아래를 자르지 위를 세우지 못한다.} 오늘 지은 둘 중 하나가 벌고 하나가
안 벌었는데, 그 차이를 \textbf{짓기 전에는 몰랐다.}

\paragraph{덤 --- 열린 갈래 하나를 닫는다} ``과제 78 시장 레코드 20건 사람
태깅''이 오래 열려 있었다. \texttt{media\_push} 가 시장팝업에서 마스크
0\%인데 사람이 채워야 한다는 것이었다. 재 보니 \textbf{대응 필드가 이미
있다} --- \texttt{intervention.promotions} 가 647건 \emph{전부}에 있다.
그런데 \textbf{자유 텍스트}라 266건에 서로 다른 값이 356개이고, 개수를
축으로 쓰면 \textbf{검사 ① 에서 떨어진다}($\rho{=}{+}0.137$, $p{=}0.17$).
그리고 시장팝업 유보 104행의 검출 문턱이 0.06$\sim$0.08 인데
\texttt{mkt\_cat}(검사 넷 통과 · 유보에서도 가름)이 $+$0.026 이다.
\textbf{20건을 태깅해도 그 문턱을 넘길 자리가 아니다.} 갈래를 닫는다.

\paragraph{남은 것}
\begin{center}\small
\begin{tabular}{ll}
\hline
갈래 & 왜 \\
\hline
다음 축에 예측을 걸어라 & 학습 $\epsilon^2$ · 기준 $\rho$ \\
\texttt{anime\_medium} 은 뺄까 & 애니 $-$0.009 · 문턱 아래 \\
기준 $\rho$ 가 낮은 도메인 & 펀딩 0.25 · 아이돌 0.10 \\
왜 유보 $\epsilon^2$ 가 거꾸로인가 & 제일 이상한 수 \\
\hline
\end{tabular}
\end{center}

\paragraph{규약} \textbf{반증 조건을 미리 적으면 반증이 결과가 된다.}
``몫이 원인''은 노트 318 · 319 · 320이 세 번 확인한 법칙에서 자연스럽게
나온 가설이었고, 그래서 더 믿었다. 미리 ``안 오르면 몫이 아니다''를
적어 두지 않았으면 $-$0.0233 을 보고 다른 설명을 찾았을 것이다.
그리고 \textbf{자기 잣대를 자기에게 대라} --- 오늘 아침 순서 주장으로 노트를
물렀으면, 저녁에 표본 셋짜리 순서를 법칙이라 부르면 안 된다.

\end{document}
